PSO logró soluciones de calidad competitiva en las instancias chicas y medianas: en las
cinco compartidas con el método exacto ($n\le33$), el GAP promedio se mantuvo entre
$0{,}03\%$ y $10{,}07\%$, tocando el óptimo exacto en \texttt{E-n23-k3} y quedando a menos
de $5\%$ de él en tres de las otras cuatro. Este desempeño queda, con todo, por debajo
del reportado por \textcite{ai2009particle} para la misma representación: su SR-1,
aplicada con el algoritmo GLNPSO y un conjunto más amplio de mejoras locales, alcanza el
óptimo exacto de \texttt{A-n33-k5} (661, Tabla~1 de ese trabajo), mientras que la mejor
de las 31 corridas de este informe llegó a 697 en esa instancia; la diferencia es
atribuible a la variante de PSO más simple usada aquí
(Sección~\ref{sec:m2_limitaciones}). En las cuatro instancias agregadas para evaluar
escalabilidad ($n\ge60$, fuera del alcance del método exacto), en cambio, el GAP
promedio se ubicó entre $15{,}25\%$ y $26{,}84\%$, con el máximo en la instancia más
grande (\texttt{E-n101-k8}, $n=101$): PSO sigue entregando una solución factible y
razonable en todos los casos, pero su ventaja relativa en calidad se reduce claramente al
alejarse del rango de tamaños usado para parametrizarlo.

\subsection{Exacto vs.\ metaheurística}

La comparación de la Sección~\ref{sec:m2_comparacion} deja ver con claridad el
compromiso entre \textit{garantía} y \textit{escalabilidad} que motiva este informe. El
método exacto certifica optimalidad y, cuando puede hacerlo en tiempo razonable
(instancias con $n\le32$), es la opción preferible sin ambigüedad: no hay ninguna razón
para aceptar una solución aproximada cuando el óptimo se obtiene en menos de 10~minutos.
Pero el salto observado en el Informe~1 entre $n=32$ (576{,}1~s) y $n=33$ (17\,096{,}4~s,
sin certificar) muestra que esa tratabilidad no se degrada suavemente: se pierde de forma
abrupta con apenas un cliente adicional. PSO, en el mismo rango, no muestra ningún
quiebre comparable: su tiempo de ejecución crece de forma acotada con $n$ (de $0{,}066$~s
en $n=13$ a $0{,}306$~s en $n=101$, Tabla~\ref{tab:m2_resultados_pso}), y sigue entregando
una solución completa incluso en instancias que el método exacto nunca llegó a intentar.

El punto en que la metaheurística se vuelve preferible no es, entonces, un umbral fijo de
$n$, sino la pregunta de si la aplicación puede tolerar una solución sin garantía de
optimalidad a cambio de una respuesta en tiempo acotado. Para instancias pequeñas donde el
exacto es rápido, no hay motivo para preferir PSO; para instancias donde el exacto se
vuelve intratable (como \texttt{A-n33-k5}, y con mayor razón las cuatro instancias más
grandes de este informe, nunca intentadas por el exacto), PSO es la única alternativa de
las dos que entrega una respuesta en un tiempo práctico, aun cuando esa respuesta pueda
estar a un GAP de dos dígitos del óptimo real.

\subsection{Limitaciones y trabajo futuro}
\label{sec:m2_limitaciones}

La implementación tiene dos limitaciones identificables a partir de los propios
resultados. Ninguna se relaciona con la parametrización en sí, que se realizó de forma
completa y sistemática con Optuna, sobre un espacio de búsqueda y un presupuesto de
evaluaciones explícitos (Sección~\ref{sec:parametrizacion}); ambas provienen de
decisiones de diseño del algoritmo y del protocolo experimental. Primero, se usó PSO
canónico (gbest global, dos términos de velocidad, Ecuación~\ref{eq:m2_velocidad}),
no la variante GLNPSO de \textcite{ai2009particle}, que agrega términos de mejor local
(\textit{lbest}) y de mejor vecino según razón aptitud-distancia (\textit{nbest}); esta
simplificación probablemente explica en parte por qué el GAP obtenido aquí es algo mayor
al reportado por los autores originales para instancias de tamaño comparable. Segundo,
el número de iteraciones $T$ se deriva de un presupuesto fijo de evaluaciones igual para
todas las instancias ($T=\lfloor 50\,000/N\rfloor$, Sección~\ref{sec:parametrizacion}),
sin ajustarse a que la dimensión del espacio de búsqueda ($n+2m$) crece con el tamaño de
la instancia; el análisis de convergencia (Sección~\ref{sec:m2_comparacion}) mostró que
la curva mediana de \texttt{E-n101-k8} todavía descendía en la última iteración
disponible, a diferencia de \texttt{A-n32-k5}, que ya se había estabilizado bastante
antes de agotar su presupuesto de iteraciones. A esto se suma que las instancias de
tuning utilizadas (hasta $n=76$, Sección~\ref{sec:parametrizacion}) no cubren todo el
rango evaluado (hasta $n=101$); esto no es una falla del proceso de parametrización, que
se ejecutó igual para todas las escalas, sino una propiedad esperable de cualquier
calibración empírica sobre un conjunto finito de instancias: los parámetros encontrados
generalizan mejor cerca del rango en que fueron ajustados. Ambos factores (el presupuesto
de iteraciones fijo y el alcance del tuning) son consistentes con el mayor GAP observado
en las instancias más grandes, y ayudan a explicarlo (Tabla~\ref{tab:m2_resultados_pso}).

Una limitación potencial de la propia representación SR-1, señalada por
\textcite{ai2009particle} (Tabla~2 de ese trabajo, instancia \texttt{vrpnc9}), es que el
mecanismo de asignación de clientes a un número fijo de vehículos puede, en casos raros,
dejar clientes sin asignar. En las $279$ corridas realizadas para este informe (31
repeticiones $\times$ 9 instancias) esto no ocurrió en ningún caso, por lo que no fue una
limitación práctica aquí, pero se documenta como un riesgo inherente al esquema de
codificación, no eliminado por la implementación.

De estas limitaciones se desprenden mejoras concretas para trabajo futuro: incorporar
los términos de vecindario social de GLNPSO (\textit{lbest}/\textit{nbest}) para acercar
el desempeño al reportado en la literatura de referencia; escalar el número de
iteraciones con la dimensión del problema en vez de mantenerlo fijo por presupuesto de
evaluaciones; ampliar el conjunto de instancias de tuning para cubrir el rango completo
de tamaños evaluados; e hibridar la búsqueda con procedimientos de mejora local
adicionales más allá del 2-opt ya embebido en la decodificación (p. ej.\ \textit{or-opt}
o intercambios entre rutas), en la línea de lo discutido para PSO memético en el
Capítulo~3.
